\documentclass[a4paper,french,12pt, titlepage]{article}

%%%% paquetages %%%%
\usepackage[french,english]{babel}  %définition de la langue
\usepackage[T1]{fontenc}
\usepackage[utf8]{inputenc} %définition de l'encodage
\usepackage{fullpage} %pour réduire les marges
\usepackage{graphicx} %pour figures
\usepackage{comment}
\usepackage{xcolor}
\usepackage{listings}
\usepackage{amsmath}
\usepackage{tikz}
\usepackage{hyperref}
\usepackage{wrapfig}
\usepackage{fancyhdr} %headers
\usepackage{glossaries}
\usepackage{todonotes}
\usepackage[export]{adjustbox}

\usepackage{titling}

\makeglossaries

\newglossaryentry{latex}
{
    name=latex,
    description={Is a markup language specially suited 
    for scientific documents}
}

\newglossaryentry{maths}
{
    name=mathematics,
    description={Mathematics is what mathematicians do}
}

\renewcommand{\maketitle}{%
  \maketitlehooka
  \maketitlehookb
  \maketitlehookc
  \maketitlehookd
}

\renewcommand{\maketitlehooka}{% 
    \vspace{-5.5cm}
    
    \noindent
    \raisebox{15ex}{\includegraphics[height=10ex]{logos/logo_polytech_full.png}}
    \hfill\raisebox{16ex}{\includegraphics[height=10ex]{logos/logo_lig_new.png.png}}
    \space\space\space
    \raisebox{12ex}{\includegraphics[height=15ex]{logos/logo_inp_uga.png}}
    
    \vfill
    
    \bigskip
    \begin{center} \large
    Alexandre Lithaud
    
    INFO4 - Polytech Grenoble
    
    Rapport de stage 2022/2023
    
    \vfill
    \begin{Large}
    \textbf{Contribution au projet NixOS Compose}
    \end{Large}
    
    \vfill
    
    Tome Principal
    
    \begin{small}
    ET
    \end{small}
    
    Annexe
    
    \end{center}
}

\renewcommand{\maketitlehookd}{% 
    \vfill{}  \large\par\noindent  
    \begin{center}
    2022/2023\\
    17 Avril 2023 - 28 Juillet 2023
    \end{center}
    \vspace{-0.5cm}
}


\newcommand{\makefooter}{%
  \makefooterhooka
}

\newcommand{\makefooterhooka}{% 
    \begin{center}
        \begin{Large}
        DOS DU RAPPORT
        \end{Large}
    \end{center}
    
    
    %\vspace{-5.5cm}
    %\noindent
    \textbf{Etudiant} : Alexandre Lithaud
    \hfill \textbf{Année d’étude dans la spécialité :}
    
    \hfill INFO4 2022/2023
    
    \hfill
    
    \textbf{Entreprise} : Laboratoire d'informatique de Grenoble 

    \textbf{Adresse complète} : Bâtiment IMAG, 700, AV. Centrale, 38401
Saint Martin d'Hères

    \textbf{Téléphone (standard)} : 07.87.30.90.36
    
    \hfill
    
    \textbf{Responsable administratif} : Noel de Palma 

    \textbf{Téléphone} : 04.57.42.14.78

    \textbf{Courriel} : noel.de-palma@univ-grenoble-alpes.fr

    \hfill
    
    \textbf{Tuteur de stage (organisme d’accueil)} : Olivier Richard

    \textbf{Téléphone} : 06.32.29.09.18

    \textbf{Courriel} : olivier.richard@imag.fr
    
    \hfill
    
    \textbf{Enseignant-référent} : Nicolas Palix

    \textbf{Téléphone} : 04.57.42.15.38 

    \textbf{Courriel} : nicolas.palix@imag.fr 

    \hfill
    
    \textbf{Titre} : Contribution au projet NixOS Compose
    
    \hfill

    \textbf{Résumé} : En 4ème année d'ingénieur en informatique, j'ai eu
l'opportunité de faire un stage de 15 semaines au Laboratoire
Informatique de Grenoble (LIG), au sein de l'équipe DATAMOVE.\newline

Durant ce stage, j'ai eu comme objectif d'utiliser et d'améliorer
l'outil NixOS-Compose, ainsi que de créer différentes compositions dans
l'optique de les utiliser à une fin de recherche. NixOS-Compose (ou NXC)
est un logiciel créé par l'équipe, permettant de décrire une
infrastructure complexe de plusieurs nœuds, en mettant l'accent sur la
reproductibilité et la simplicité de mise en place. De plus, j'ai été
amené à contribuer à la maintenance de logiciels tels que OAR et EAR,
améliorant leur stabilité et performances par le biais de mise à jour.
Le tout en utilisant le système Grid5000 qui m'a permis de tester mes
développements dans un environnement réel.\newline

Durant ce rapport, vous allez suivre la création des différentes
compositions que j'ai créée dans le but de tester les performances de
plusieurs systèmes de fichiers distribués dans le réseau de nœud
Grid5000.
}

\title{\Huge \bfseries\begin{center}Contribution au projet NixOS
Compose\end{center}}
\author{Alexandre Lithaud}
\date{17 Avril 2023 - 28 Juillet 2023}

\lstset{language=C++,
                basicstyle=\ttfamily,
                keywordstyle=\color{blue}\ttfamily,
                stringstyle=\color{red}\ttfamily,
                showstringspaces=false,
                %commentstyle=\color{magenta}\ttfamily,
                morecomment=[l][\color{magenta}]{\#}
}

\newcommand*\xor{\mathbin{\oplus}}
\newcommand{\paragraphnewline}[1]{\hypertarget{par#1}{\paragraph{#1}\mbox{}}}


\pagestyle{fancy}
\fancyhf{}
\rhead{Alexandre Lithaud - 2022/2023}
\lhead{Rapport de stage}
\rfoot{Page \thepage}
\renewcommand{\headrulewidth}{1pt}
\renewcommand{\footrulewidth}{1pt}
\setlength{\headheight}{15pt}
\headsep = 1.0cm

% references

\usepackage{csquotes}
\usepackage{biblatex}
\bibliography{references.bib}


% \makenoidxglossaries

\begin{document}
\selectlanguage{french}

\begin{titlingpage}
\maketitle
\end{titlingpage}

\begin{center}
    \item \paragraphnewline{Remerciements}
\end{center}

Je tiens à tout d'abord à remercier le Laboratoire Informatique de
Grenoble et tous ces membres pour l'accueil chaleureux que j'ai reçu à
mon arrivée au laboratoire ainsi que pour l'ambiance générale du stage
qui a été exemplaire.\newline

Je remercie également Monsieur Olivier Richard et Monsieur Nicolas
Palix, respectivement mon tuteur et mon référent de stage pour leurs
conseils ainsi que leurs pédagogies qui m'ont permis de réaliser mes
missions dans les meilleures conditions possibles et de grandement
monter en compétence durant ce stage.\newline

Enfin, je suis reconnaissant envers Quentin GUILLOTEAU et Adrien FAURE,
respectivement doctorant et chercheur au Laboratoire Informatique de
Grenoble pour les inestimables conseils et les réponses dispensés lors
de mes différentes missions.

\newpage

\selectlanguage{french}
\begin{center}
    \item \paragraphnewline{Résumé}
\end{center}

En 4ème année d'ingénieur en informatique, j'ai eu l'opportunité de
faire un stage de 15 semaines au Laboratoire Informatique de Grenoble
(LIG), au sein de l'équipe DATAMOVE.\newline

Durant ce stage, j'ai eu comme objectif d'utiliser et d'améliorer
l'outil NixOS-Compose, ainsi que de créer différentes compositions dans
l'optique de les utiliser à une fin de recherche. NixOS-Compose (ou NXC)
est un logiciel créé par l'équipe, permettant de décrire une
infrastructure complexe de plusieurs nœuds, en mettant l'accent sur la
reproductibilité et la simplicité de mise en place. De plus, j'ai été
amené à contribuer à la maintenance de logiciels tels que OAR et EAR,
améliorant leur stabilité et performances par le biais de mise à jour.
Le tout en utilisant le système Grid5000 qui m'a permis de tester mes
développements dans un environnement réel.\newline

Durant ce rapport, vous allez suivre la création des différentes
compositions que j'ai créée dans le but de tester les performances de
plusieurs systèmes de fichiers distribués dans le réseau de nœud
Grid5000.


\textbf{\textit{mots-clés---}} Nix, Programmation Fonctionnel,
Laboratoire, NixOS, NixOS-Compose, Grid5000, Systèmes de fichiers,
Logiciel de Recherche, Maintenance, HPC, Infrastructure Distribué.

\selectlanguage{english}
\begin{center}
    \item \paragraphnewline{Abstract}
\end{center}

In my 4th year as a computer science engineer, I had the opportunity to
do a 15-week internship at the IT Laboratory of Grenoble (LIG), in the
DATAMOVE team.\newline

During this placement, my aim was to use and improve the NixOS-Compose
tool, and to create various compositions with a view to using them for
research purposes. NixOS-Compose (or NXC) is a piece of software created
by the team, enabling a complex infrastructure of several nodes to be
described, with the emphasis on reproducibility and simplicity of
implementation. I also contributed to the maintenance of software such
as OAR and EAR, improving their stability and performance through
updates. All this was done using the Grid5000 system, which enabled me
to test my developments in a real environment.\newline 

In this report, you will follow the creation of the various compositions
I created in order to test the performance of several distributed file
systems in the Grid5000 node network.

\textbf{\textit{Keywords---}} Nix, Functional Programming, Laboratory,
NixOS, NixOS-Compose, Grid5000, File Systems, Research Softwares,
Maintenance, HPC, Distributed Infrastructure

\selectlanguage{french}
\newpage

\tableofcontents
\newpage

\listoffigures

\newpage

\hypertarget{introduction}{%
\section{Introduction}\label{introduction}}

Ce rapport va représenter mon expérience de stage au Laboratoire
Informatique de Grenoble. Mon stage de 15 semaines à débuter le 17 avril
2023. Au cours de cette période j'ai eu l'opportunité de travailler sur
divers projet informatiques en lien avec les technologies de Nix, NixOS
et le HPC (\emph{High performance computing}). Ainsi que sur la
maintenance et l'amélioration de logiciel et recherche tels que OAR et
EAR. Cette opportunité m'a donné l'occasion de travailler avec le
système Grid5000, qui offre une infrastructure distribuée pour
l'exécution de travaux de recherche à grande échelle.\newline

L'objectif principal de mon stage était, en premier lieu, de contribuer
au projet au projet NixOS-Compose, un outils puissant qui facilite le
déploiement et la gestion d'environnement de développement reproductible
spécialisé pour le HPC en déployant directement plusieurs machines sur
Grid5000 à la manière de Docker Compose (d'où le nom). Afin de pourvoir
réaliser cette tache il était important de monter en compétences sur
Nix, NixOS. Grâce à cette expérience, j'ai pu approfondir ma
compréhension des principes fondamentaux de la gestion des paquets et
des environnements isolés, la configuration de système basé NixOS, le
paradigme de programmation fonctionnelle ainsi que le déploiement
d'application fonctionnelle dans un environnement d'HPC.\newline

En parallèle, j'ai participé à la maintenance et à l'amélioration de
logiciel de recherche tels que OAR et EAR en les mettant à jour avec la
dernière version de Nix par exemple. OAR joue un role crucial dans la
planification de travaux de recherche sur des infrastructures distribué
comme Grid5000 notamment. EAR quant à lui, permet d'instrumenter et donc
de quantifié les performances d'applications distribuées.J'ai pu
contribuer à l'amélioration de leur stabilité, de leurs performances et
de leurs fonctionnalités, en collaborant étroitement avec l'équipe de
développement du laboratoire.\newline

De plus, j'ai eu l'opportunité de travailler en utilisant le système
Grid5000, qui m'a permis de déployer et de tester mes compositions
directement dans un environnement réel et reproductible. Cette
expérience m'a offert une compréhension bien plus poussé sur les
méthodes de déploiement de logiciel, à l'importance de d'évolutivité et
à la gestion des ressources et à la fiabilité des systèmes
distribués.\newline

Dans ce rapport, je décrirai en détail les différentes tâches et projets
auxquels j'ai participé tout au long de mon stage, en mettant l'accent
sur les compétences acquises, les résultats obtenus et les leçons
apprises. Je présenterai également une analyse critique de mes
réalisations, ainsi que des suggestions pour des améliorations futures.
Ce rapport témoigne de ma progression en tant que professionnel de
l'informatique et des contributions significatives que j'ai apportées au
sein du LIG.\newline

\newpage

\hypertarget{contexte-du-stage}{%
\section{Contexte du stage}\label{contexte-du-stage}}

\hypertarget{le-laboratoire-informatique-de-grenoble}{%
\subsection{Le Laboratoire Informatique de
Grenoble}\label{le-laboratoire-informatique-de-grenoble}}

Test \gls{latex} test

\hypertarget{luxe9quipe-datamove}{%
\subsection{L'équipe DATAMOVE}\label{luxe9quipe-datamove}}

\hypertarget{le-projet-nixos-compose}{%
\subsection{Le Projet NixOS-Compose}\label{le-projet-nixos-compose}}

\newpage

\hypertarget{annexe}{%
\section{Annexe}\label{annexe}}

\printbibliography

\printglossaries

\begin{titlingpage}
\newpage
\makefooter
\end{titlingpage}

\end{document}
